\documentclass{book}
\usepackage{ctex}
\usepackage[cmex10]{amsmath}
\usepackage{xcolor}
\usepackage{mathtools}
\usepackage{booktabs}
\DeclareRobustCommand{\bluebold}[1]{\textbf{\textcolor[rgb]{0.31,0.63,0.86}{#1}}}

\usepackage[
    pdfpagemode=UseOutlines,
    bookmarks=true,
    bookmarksopen=true,
    bookmarksnumbered=true,
    colorlinks=true,
    linkcolor=black,
    citecolor=black,
    urlcolor=black
]{hyperref}
\usepackage{bookmark}

\raggedbottom

\begin{document}
\tableofcontents
\clearpage

管路计算

\begin{table}[htbp]
\centering
\begin{tabular}{c c c} 
\toprule
知识点 & point & 常见题型\\ [0.5ex] 
\midrule
1 & 简单管路 & 填空 \ 选择 \ 计算\\
2 & 复杂管路 & 填空 \ 选择\\
\bottomrule
\end{tabular}
\caption{Knowledge points table for pipeline calculation.}
\label{table:data}
\end{table}

\chapter{简单管路}

单根管子或由不同材质，不同管径的管子连接而成的串联管路

各段管路内流体的\bluebold{流量相同}。

\begin{equation}
q_\mathrm{v} = q_\mathrm{v1} = q_\mathrm{v2}
\end{equation}

管路的\bluebold{总阻力损失=各段管路阻力损失之和}

\begin{equation}
h_\mathrm{f,\text{总}} = h_\mathrm{f,1} + h_\mathrm{f,2}
\end{equation}

\bluebold{前面所讲的有关管路的内容，都属于简单管路}

\chapter{复杂管路}

\section{并联管路}

有分支也有汇合

主管中的流量\bluebold{等于}并联的各支管流量之和，对于不可压缩流体，则有

\begin{equation}
q_\mathrm{v} = q_\mathrm{v1} + q_\mathrm{v2}
\end{equation}

并联管路中，单位质量流体通过任何一根支管的\bluebold{擦损失相等},并基于此进行流量分配

\begin{equation}
h_\mathrm{f,1} = h_\mathrm{f,2} = h_\mathrm{f,AB}
\end{equation}

\section{分支管路}

有分支，不汇合

主管中的流量等于并联的各支管流量之和，对于不可压缩流体，则有

\begin{equation}
q_\mathrm{v} = q_\mathrm{vA} + q_\mathrm{vB}
\end{equation}

单位质量流体在各支管流动终了时的总机械能与能量损失之和相等

\begin{equation}
gz_\mathrm{C} + \frac{u^2_c}{2} + \frac{p_c}{\rho} = gz_\mathrm{A} + \frac{u^2_A}{2} + \frac{p_A}{\rho} + h_\mathrm{f, CA} = gz_\mathrm{B} + \frac{u^2_B}{2} + \frac{p_B}{\rho} + h_\mathrm{f, CB}
\end{equation}

\end{document}
